% ======================================================================
% SECTION 2: METHODS
% Six thematic subsections: instruction inputs, robot 2030 NeuroSpeed,
% sensors and sample table, xyz coordinate mapping, iterations plus
% LLM tournament, code execution environment.
% Bracketed prompts for future Claude Code Opus 4.7 1M Max processing.
% Do not process the bracketed prompts here.
% ======================================================================

\subsection{Instruction Inputs}
\label{sec:methods-inputs}

[METHODS SUBSECTION 1 PROMPT - PROCESS DURING THE PAPER POPULATION PASS]

[Document the 12 hand-authored instruction files that the
2030-gbm-1min repository was generated from. Read the directory
kevinkawchak/physical-ai-oncology-trials/competitions/instructions/one_minute_variant/
and enumerate every file by exact name and one-line summary: README.md
(scope and per-commit roadmap), commit_01_overview_1min.md (project
narrative and per-commit roadmap), commit_02_sensors_1min.md (sensor
specification at 200 channels mixed 1 kHz plus 10 kHz), commit_03_xyz_4arm.md
(deterministic sensor to xyz transformation across 4 cooperating arms),
commit_04_iterations_1min.md (16-iteration deterministic sweep design
and YAML configuration), commit_05_competition_1min.md (LLM tournament
and composite score formula), file_size_pyramid_1min.md (Git commit
caps and Zenodo deposition strategy), glioblastoma_context_1min.md
(IDH-wildtype glioblastoma clinical context), multi_arm_coordination.md
(1 kHz heartbeat broadcast bus protocol), robot_specification_neurospeed.md
(2030 Medtronic NeuroSpeed 1.0 7-DOF DH parameters and per-arm tool
assignment), sensor_specification_10khz.md (50 channels per arm at
1 kHz commands plus 10 kHz force), zenodo_archive_protocol.md
(Zenodo L0 raw deposition protocol at DOI 10.5281/zenodo.18445179).
State explicitly that these 12 files together amount to roughly 130 KB
of read-only context, and that the entire 2030-gbm-1min/ repository
was generated by Claude Code Opus 4.7 1M Max in a single sequence of
seven commits read against these instructions. State the
approximations made (number of files, file sizes, sampling rates,
iteration count) and which were left for the downstream pass; cross
reference Limitations and Future Work section
(\S\ref{sec:limitations}) for the detailed accounting.]

\subsection{2030 Medtronic NeuroSpeed 1.0 Robot}
\label{sec:methods-robot}

[METHODS SUBSECTION 2 PROMPT - PROCESS DURING THE PAPER POPULATION PASS]

[Read
kevinkawchak/physical-ai-oncology-trials/competitions/instructions/one_minute_variant/robot_specification_neurospeed.md
and
kevinkawchak/robotic-surgeries/2030-gbm-1min/config/kinematics_4arm.yaml
in full. Reproduce the 7-DOF DH parameter table per arm using the
column type {>{\raggedright\arraybackslash}p{Xcm}} for each column
(prepend \raggedright\arraybackslash to every width value). Include
joint angular velocity caps, tip force limits (per-arm 5.0 N tip, 1.0 N
lateral), cumulative 4-arm tip force cap (12 N on the patient frame),
5 ms E-stop budget, 100 microsecond emergency arm-park trigger,
inter-arm clearance (8 mm minimum), and 1 kHz heartbeat broadcast bus
protocol (32-byte frame, 1 ms deadline, 3 ms watchdog). State
explicitly the per-arm tool assignment: arm 1 hybrid ultrasonic plus
waterjet plus pulsed plasma, arm 2 bipolar coagulation plus
irrigation, arm 3 suction plus tissue collection, arm 4 0.5 T iMRI
plus 5-ALA fluorescence plus ultrasound. Cite iec-80601-2-77,
iec-62304, and rosa-one-brain. Reference the 4-phase 60-second
timeline (dural opening final 5 s, bulk resection 40 s, margin
assessment 10 s, hemostasis withdrawal 5 s) read from
commit_01_overview_1min.md, and the per-arm contribution targets read
from commit_04_iterations_1min.md.]

\subsection{Sensor Specification (54 columns by 1001 rows)}
\label{sec:methods-sensors}

[METHODS SUBSECTION 3 PROMPT - PROCESS DURING THE PAPER POPULATION PASS]

[Read
kevinkawchak/physical-ai-oncology-trials/competitions/instructions/one_minute_variant/sensor_specification_10khz.md,
kevinkawchak/robotic-surgeries/2030-gbm-1min/schemas/sensor_record_4arm.schema.json,
kevinkawchak/robotic-surgeries/2030-gbm-1min/schemas/sensor_record_4arm.proto,
and
kevinkawchak/robotic-surgeries/2030-gbm-1min/schemas/sensor_record_4arm.avsc
in full. Document the 200-channel sensor stack: 50 channels per arm
x 4 arms = 200 channels at mixed 1 kHz commands plus 10 kHz force.
Highlight the exceptional processing ability that no human team
could produce in due time. Specifically reference the file
kevinkawchak/robotic-surgeries/2030-gbm-1min/outputs/sensors/sensor_sample_4arm.csv
as a 54 columns by 1001 rows table (the CSV holds 50 sensor channels
per arm collapsed plus 4 metadata columns spanning timestamp_us,
arm_id, phase, command_state, across 1001 time rows; the JSONL
companion at outputs/sensors/sensor_sample_4arm.jsonl carries the
full per-arm record). Render a small example slice as a LaTeX table
using {>{\raggedright\arraybackslash}p{Xcm}} columns. Reference the
release-aggregate L0 raw archive that lives on Zenodo only (cite
zenodo) at DOI 10.5281/zenodo.18445179 (cite repo-robotic-surgeries).
State the file size cap policy from
kevinkawchak/physical-ai-oncology-trials/competitions/instructions/one_minute_variant/file_size_pyramid_1min.md
and the 1 KB pointer JSON pattern at
kevinkawchak/robotic-surgeries/2030-gbm-1min/data/sensor_l0_raw_4arm.zenodo_pointer.json.
Cite apache-arrow for the Parquet format used by the L1 to L3
aggregates.]

\subsection{Per-Arm xyz Coordinate Mapping}
\label{sec:methods-xyz}

[METHODS SUBSECTION 4 PROMPT - PROCESS DURING THE PAPER POPULATION PASS]

[Read
kevinkawchak/physical-ai-oncology-trials/competitions/instructions/one_minute_variant/commit_03_xyz_4arm.md,
kevinkawchak/physical-ai-oncology-trials/competitions/instructions/one_minute_variant/multi_arm_coordination.md,
kevinkawchak/robotic-surgeries/2030-gbm-1min/src/mapping/sensor_to_xyz_4arm.py,
kevinkawchak/robotic-surgeries/2030-gbm-1min/src/control/robot_loop_4arm.cpp,
kevinkawchak/robotic-surgeries/2030-gbm-1min/src/coordination/arm_heartbeat.cpp,
and
kevinkawchak/robotic-surgeries/2030-gbm-1min/schemas/xyz_command_4arm.schema.json
in full. Document the deterministic sensor to xyz transformation
across 4 cooperating arms: forward kinematics from 7-DOF joint state
to end effector position in patient frame, inverse kinematics with
tolerance configurable in {1e-6, 1e-5, 1e-4, 1e-3}, and the
1 kHz heartbeat coordination layer with 32-byte frames. Reference the
per-arm xyz command samples at
kevinkawchak/robotic-surgeries/2030-gbm-1min/outputs/xyz_mapping/xyz_trace_sample_arm1.csv
through xyz_trace_sample_arm4.csv (4 files of 60 rows each) and the
ASCII per-second xyz path overlay at outputs/xyz_mapping/xyz_path_4arm.txt.
Render the per-arm command state distribution as a LaTeX table using
{>{\raggedright\arraybackslash}p{Xcm}} for each column. Cite
iec-80601-2-77 for the force envelope and iec-62304 for the
safety-critical software lifecycle. Highlight that all 240 commands
in the 4-file sample resolve to command_state=EMIT (no rejections,
no E-stop triggers in the sample window).]

\subsection{16-Iteration Sweep and On-Premises LLM Tournament}
\label{sec:methods-iterations}

[METHODS SUBSECTION 5 PROMPT - PROCESS DURING THE PAPER POPULATION PASS]

[Read
kevinkawchak/physical-ai-oncology-trials/competitions/instructions/one_minute_variant/commit_04_iterations_1min.md,
kevinkawchak/physical-ai-oncology-trials/competitions/instructions/one_minute_variant/commit_05_competition_1min.md,
kevinkawchak/robotic-surgeries/2030-gbm-1min/config/iterations.yaml,
kevinkawchak/robotic-surgeries/2030-gbm-1min/src/simulation/iterate_1min.py,
kevinkawchak/robotic-surgeries/2030-gbm-1min/src/simulation/runner_1min.rs,
kevinkawchak/robotic-surgeries/2030-gbm-1min/src/metrics/compute_1min.py,
kevinkawchak/robotic-surgeries/2030-gbm-1min/src/llm/compare_agent_1min.py,
kevinkawchak/robotic-surgeries/2030-gbm-1min/prompts/comparison_prompt_1min.md,
and
kevinkawchak/robotic-surgeries/2030-gbm-1min/docs/comparison_methodology.md
in full. Document the 16-iteration deterministic sweep across noise
(0.01 to 0.05 mm sensor noise sigma), gain (0.8 to 1.2 force feedback
gain), inverse kinematics tolerance (1e-6 to 1e-3), and heartbeat
jitter (0 to 50 microseconds). State the deterministic seed 20260510
and the bit-deterministic property for fixed seed plus parameter
tuple. Document the composite score formula under the frozen weighted
formula: quality 0.40 + time 0.25 + cost 0.20 + safety 0.10 +
patient_experience 0.05. State that the on-premises LLM tournament
runs a 4-entity bracket (robot vs robot default, mixed 2 robot plus
2 human optional) with Anthropic claude-opus-4-7 as the default judge
and Ollama as the optional on-premises backend (cite claude-opus-47,
claude-code, ollama, vllm). Reference the per-iteration L1 50 ms,
L2 1 s, L3 phase, and events Parquet aggregates emitted to
kevinkawchak/robotic-surgeries/2030-gbm-1min/outputs/iterations/
and the cross-iteration index.jsonl manifest plus DuckDB analytical
store (cite duckdb). Note the structural-time-dimension caveat that
is preserved in every LLM judge rationale (1-minute robot vs 1-hour
human baseline).]

\subsection{Code Execution Environment}
\label{sec:methods-execution}

[METHODS SUBSECTION 6 PROMPT - PROCESS DURING THE PAPER POPULATION PASS]

[Read
kevinkawchak/robotic-surgeries/2030-gbm-1min/README.md (Quick Start
section, lines covering Linux, MacOS, Windows, NVIDIA A100, and
Claude Code recipes),
kevinkawchak/robotic-surgeries/2030-gbm-1min/pyproject.toml,
kevinkawchak/robotic-surgeries/2030-gbm-1min/docker-compose.yml,
kevinkawchak/robotic-surgeries/2030-gbm-1min/src/simulation/Cargo.toml,
and
kevinkawchak/robotic-surgeries/.github/workflows/ci.yml
in full. Document the cross-platform setup recipes for Linux
(Ubuntu 22.04 LTS or later), MacOS (M3 Ultra or Intel), Windows 11
with PowerShell 7, the NVIDIA A100 GPU recipe with optional CUDA
feature flag, and Claude Code (CLI, web, or IDE plugin) execution.
Render the expected single-iteration wall-clock per platform (M3 Ultra
26 to 32 s, A100 12 s) as a LaTeX table using
{>{\raggedright\arraybackslash}p{Xcm}} for each column. Reference the
ruff format, ruff check, yamllint, and file size cap CI gates that
run on Python 3.10, 3.11, and 3.12. State the verification recipe
from 2030-gbm-1min/README.md. Note that the C++20 control loop and
C++20 1 kHz heartbeat layer compile on Linux (g++), MacOS (clang++),
and Windows (MSVC). Cite apache-arrow for the Parquet output format.]
